\documentclass[12pt]{article}
\usepackage{parskip}

\author{Andrea Riolo Vinciguerra}
\title{Compilation Techniques: Laboratory Project\\ Project Report}
\date{\today}

\begin{document}

\maketitle
\tableofcontents
\clearpage

\section{Introduction}

This project is organized in a single codebase that houses two compilers:
MiniImp and MiniFun. While both share common background code, the core logic is
kept strictly separate from the main user features. This setup allows users to
easily choose between generating a control flow graph, running an interactive
interpreter, or executing a code file.

A main focus of the design is modularity. By keeping the code well-organized and
adaptable, the system is much easier to maintain and expand over time. Comments
are included (more or less helpful) throughout the source files to make
comprehension easier.

The code is organized by separating distinct functional blocks into different
files. For example, the core evaluator (or interpreter) and the Abstract Syntax
Tree definitions for both compilers are stored in a file called
\texttt{engine.ts}. The rest of the functionality builds upon this file. For
example, MiniImp extends this core engine to construct its control flow graph
and handle optimizations, while MiniFun supports a complete polymorphic type
system.

I chose TypeScript as the programming language for two main reasons. First, I
wanted the opportunity to master a language I was less familiar with. Second,
TypeScript is excellent at handling union types. This feature (used extensively
throughout the program) makes writing compiler code much simpler by avoiding the
bulky, complex structures often needed in languages like C++ or Rust.

Because Node.js works smoothly across different operating systems, the project
can be packaged into a single, easy-to-run executable file for both macOS and
Linux. On Windows, however, the operating system does not support shebangs (the
special line of code that automatically tells the system how to run a script).
As a result, Windows users need to run the program manually by invoking the
\texttt{node} command.

\section{Running the code}

The software can be run in two ways: automatically or manually. The automatic
method is designed specifically for Linux or macOS. Instead, the manual approach
works on any operating system.

Either approach requires that Node.js and its package manager, \texttt{npm}, are
installed and properly configured in the system's \texttt{PATH}.

\subsection{Automatic run}

Inside the main project folder, there is a helpful script called
\texttt{run.sh}. This script handles the setup automatically by downloading the
required files, building the chosen compiler (MiniImp or MiniFun), and launching
the program.

Running either compiler is straightforward. For example, to start MiniImp, this
command is sufficient:
\begin{verbatim}
$ ./run.sh miniimp --help
\end{verbatim}

Similarly, for MiniFun:
\begin{verbatim}
$ ./run.sh minifun --help
\end{verbatim}

Any extra instructions or arguments should be typed directly after the
compiler's name. As shown in the examples above, adding the \texttt{--help} flag
is a great way to see a full list of all the available commands.

\subsection{Manual compilation}

Manual compilation is available. In such case, these commands must be executed
upon downloading the repository:
\begin{verbatim}
$ npm ci
$ npm run build:miniimp     # or npm run build:minifun
$ dist/miniimp --help       # or dist/minifun --help
\end{verbatim}

For Windows, the last step involves invoking \texttt{node} explicitly, like so:
\begin{verbatim}
$ node dist/miniimp --help  # or node dist/minifun --help
\end{verbatim}

\section{MiniFun: Functional Language}

MiniFun is a functional language supporting three types of first-class values:
closures (functions that store the environment at the time of definition),
integers and booleans. It supports let bindings and recursion.

%block graph has 'merge' block to handle the case where, after an if, there's
%the end of the program. without 'merge' block, we would have an empty linear
%block (fake skip is not a command to display). while it has been done mainly
%for debugging purposes, merge and linear blocks are fundamentally equivalent to
%each other while encoding the difference needed for SSA-based algorithms.

\end{document}
